% 附录 C 量子物理与原子物理（Quantum and atomic physics）
\chapter{量子物理与原子物理}

\section{量子力学}

本节简要回顾量子力学的若干基本特征。量子力学中的基本对象是波函数（又称态函数）
$\psi$——它包含从观察获得的关于系统的全部知识。波函数 $\psi$ 是复函数，
$|\psi|^{2}=\psi\psi^{*}$ 是概率密度。\fn{1}{此处符号 $*$ 表示复共轭。}每个可观察
量都对应一个厄米算符。进行测量时，得到的结果是该算符的本征值之一。厄米算符具有
实本征值与彼此正交的本征函数。若算符 $\hat{A}$（算符都戴“帽子”）的第 $i$ 个
本征函数为 $\phi_i$（设已归一化）、本征值为 $a_i$，则
\begin{equation*}
\hat{A}\phi_i = a_i\phi_i.\tag{C.1}
\end{equation*}
测量算符 $\hat{A}$ 之后得到的期望值由
\begin{equation*}
\langle\hat{A}\rangle = \int\mathrm{d}\tau\,\psi^{*}\hat{A}\psi,\tag{C.2}
\end{equation*}
给出，$\mathrm{d}\tau$ 是体积元。波函数 $\psi$ 可以用 $\hat{A}$ 的本征函数展开：
\begin{equation*}
\psi = \sum_i c_i\phi_i,\tag{C.3}
\end{equation*}
利用 $\phi_i$ 的正交归一性，$\langle\hat{A}\rangle$ 于是为
\begin{equation*}
\langle\hat{A}\rangle = \sum_i|c_i|^{2}a_i,\tag{C.4}
\end{equation*}
得到第 $i$ 个本征值 $a_i$ 的概率是 $|c_i|^{2}$。测量是一种剧烈的过程：取波函数
$\psi$，测量与算符 $\hat{A}$ 相联系的物理量并得到结果 $a_n$（$\hat{A}$ 的第 $n$
个本征值），系统就被迫进入态函数 $\phi_n$——至少哥本哈根诠释如此。

两个算符 $\hat{A}$ 与 $\hat{B}$ 的对易子定义为
$[\hat{A},\hat{B}]=\hat{A}\hat{B}-\hat{B}\hat{A}$。定义对易子的表达式称为对易
关系。

\begin{example}{C.1}
对易子 $[\hat{x},\hat{p}]=\mathrm{i}\hbar$，其中 $\hat{x}$ 与 $\hat{p}$ 分别是
位置与动量算符。
\end{example}

若两个算符 $\hat{A}$ 与 $\hat{B}$ 对易（$[\hat{A},\hat{B}]=0$），称它们相容
（compatible）：测量其一不以任何方式影响另一者的取值。若不对易，则二者之间存在
不确定关系：
\begin{equation*}
\Delta A\,\Delta B \geq \frac{1}{2}|\langle\mathrm{i}[\hat{A},\hat{B}]\rangle|,\tag{C.5}
\end{equation*}
$\Delta A\equiv\sqrt{\langle(\hat{A}-\langle\hat{A}\rangle)^{2}\rangle}$。

不做测量时 $\psi$ 的时间依赖由薛定谔方程给出：
\begin{equation*}
\hat{\mathcal{H}}\psi = \mathrm{i}\hbar\frac{\mathrm{d}\psi}{\mathrm{d}t}\tag{C.6}
\end{equation*}
$\hat{\mathcal{H}}$ 是哈密顿量。算符期望值的时间依赖为
\begin{equation*}
\frac{\mathrm{d}\langle\hat{A}\rangle}{\mathrm{d}t} = \frac{1}{\mathrm{i}\hbar}\langle[\hat{A},\hat{\mathcal{H}}]\rangle,\tag{C.7}
\end{equation*}
前提是算符 $\hat{A}$ 自身不依赖时间。与哈密顿量对易的可观察量是守恒量，称为运动
恒量（constant of the motion）或好量子数。

\section{狄拉克左右矢记号}

本节回顾狄拉克（Paul Dirac）发明、本书偶有使用而磁学文献中常见的一套记号。矢量
$\mathbf{a}$ 与 $\mathbf{b}$ 的标量积为 $\mathbf{a}\cdot\mathbf{b}=\sum_ia_ib_i$。
$\mathbf{a}$ 的长度为 $\sqrt{\mathbf{a}\cdot\mathbf{a}}$，必须为正实数。这一记号
使内积在 $\mathbf{a}$ 与 $\mathbf{b}$ 之间显得完全对称，但其中藏有微妙之处。若
矢量是复的，标量积应写为
\begin{equation*}
\mathbf{a}^{\dagger}\cdot\mathbf{b} = \sum_ia_i^{*}b_i = (a_1^{*}\quad a_2^{*}\quad a_3^{*})\begin{pmatrix}b_1\\b_2\\b_3\end{pmatrix}.\tag{C.8}
\end{equation*}
这一定义保证 $a^{\dagger}\cdot a=\sum_i|a_i|^{2}$ 为正实数，复矢量的长度才有良好
定义。式 C.8 强调：标量积中的两个矢量其实不应等量齐观——一个是行矢量、另一个是
列矢量；一个取了复共轭、另一个没有。数学家说这两个矢量其实“住”在不同的空间里：
b 住在矢量空间，a 住在它的对偶空间（dual space）。要把矢量变成其对偶，须取伴随
（adjoint），以 † 号标记（意思是“取复共轭并转置”）：
\begin{equation*}
\begin{pmatrix}a_1\\a_2\\a_3\end{pmatrix}^{\dagger} = (a_1^{*}\quad a_2^{*}\quad a_3^{*}).\tag{C.9}
\end{equation*}
然而在量子力学中，态函数有时是矢量（如 $\mathbf{a}=\begin{pmatrix}a_1\\a_2\end{pmatrix}$），
有时是连续函数 $a(x)$。若是函数，则 $a(x)$ 与 $b(x)$ 的标量积为
\begin{equation*}
a^{\dagger}\cdot b = \int\mathrm{d}x\,a^{*}(x)b(x).\tag{C.10}
\end{equation*}
狄拉克用把态函数写成右矢（ket）$|a\rangle$ 的办法绕开了这一记号难题——无论处理
的是有限分量矢量还是函数（后者可视为无穷维矢量）。这只是记号：$|a\rangle+|b\rangle=|c\rangle$
就是 $\mathbf{a}+\mathbf{b}=\mathbf{c}$（或等价地 $\overrightarrow{a}+\overrightarrow{b}=\overrightarrow{c}$
或 $a(x)+b(x)=c(x)$）的另一种写法。可以用伴随把右矢 $|a\rangle$ 变成左矢（bra）
$\langle a|$：
\begin{equation*}
|a\rangle^{\dagger} = \langle a|.\tag{C.11}
\end{equation*}
$\langle a|$ 与 $|b\rangle$ 的标量积写作
\begin{equation*}
\langle a|b\rangle = \sum_ia_i^{*}b_i = \int\mathrm{d}x\,a^{*}(x)b(x),\tag{C.12}
\end{equation*}
即 bra-(c)-ket——“括号”（狄拉克在此展现出幽默感的证据）。算符期望值的表达式
可写为
\begin{equation*}
\langle\hat{A}\rangle = \langle\psi|\hat{A}|\psi\rangle.\tag{C.13}
\end{equation*}
这一记号传达了标量积的不对称性，被广泛使用并大大简化了许多计算。它在标记特定
态时尤其有用，例如：
\begin{equation*}
|\text{薛定谔的猫}\rangle = \frac{|\text{活}\rangle+|\text{死}\rangle}{\sqrt{2}}.\tag{C.14}
\end{equation*}

\section{玻尔模型}

玻尔模型是描述氢原子中电子运动的半经典模型：它给出正确的结果，却是出于错误的
理由。尽管如此，它足够简单，可用于快速推导诸如原子轨道半径对核电荷 Z 与有效
质量的依赖之类结果。

考虑电子绕电荷 +Ze 的原子核运动（见图 C.1）。令静电力与向心力相等：
\begin{equation*}
\frac{Ze^{2}}{4\pi\epsilon_0 r^{2}} = \frac{m_{\mathrm{e}}v^{2}}{r}.\tag{C.15}
\end{equation*}
电子角动量 $m_{\mathrm{e}}vr$ 以 $\hbar$ 为单位量子化：
\begin{equation*}
m_{\mathrm{e}}vr = n\hbar,\tag{C.16}
\end{equation*}
n 是整数（主量子数）。于是轨道半径可写为
\begin{equation*}
r = \frac{4\pi\epsilon_0\hbar^{2}n^{2}}{Ze^{2}m_{\mathrm{e}}} = \left(\frac{4\pi\epsilon_0\hbar^{2}}{m_{\mathrm{e}}e^{2}}\right)\frac{n^{2}}{Z} = \frac{a_0n^{2}}{Z},\tag{C.17}
\end{equation*}
\mnote{这些结果可写作 $a_0=\frac{h}{m\alpha c}$ 与 $E=-\frac{1}{2}m(\alpha c)^{2}\frac{Z^{2}}{n^{2}}$，$\alpha$ 为精细结构常数。}
$a_0$ 是玻尔半径：
\begin{equation*}
a_0 = \frac{4\pi\epsilon_0\hbar^{2}}{m_{\mathrm{e}}e^{2}}.\tag{C.18}
\end{equation*}
能量 E 为
\begin{align*}
E &= \frac{1}{2}mv^{2} - \frac{Ze^{2}}{4\pi\epsilon_0 r}\tag{C.19}\\
&= -\frac{1}{2}\left(\frac{Ze^{2}}{4\pi\epsilon_0 r}\right)\tag{C.20}\\
&= -\frac{\hbar^{2}}{2m_{\mathrm{e}}}\frac{1}{a_0^{2}}\frac{Z^{2}}{n^{2}}\tag{C.21}\\
&= -\frac{m_{\mathrm{e}}e^{4}}{(4\pi\epsilon_0\hbar)^{2}}\frac{Z^{2}}{n^{2}}.\tag{C.22}
\end{align*}

\begin{figure}
\centering
\includegraphics[width=0.35\textwidth]{figC_01.pdf}
\caption{玻尔模型。}
\end{figure}

\section{轨道角动量}

角动量算符为
\begin{equation*}
\hbar\hat{\mathbf{L}} = \hat{\mathbf{r}}\times\hat{\mathbf{p}} = -\mathrm{i}\hbar\,\hat{\mathbf{r}}\times\nabla.\tag{C.23}
\end{equation*}
角动量的 z 分量的算符为
\begin{align*}
\hbar\hat{L}_z = \mathrm{i}\hbar\left[y\frac{\partial}{\partial x} - x\frac{\partial}{\partial y}\right]\tag{C.24}\\
= -\mathrm{i}\hbar\frac{\partial}{\partial\phi}\tag{C.25}
\end{align*}
其本征函数为 $e^{\mathrm{i}m_l\phi}$、本征值 $m_l\hbar$，$m_l$ 是磁量子数。下列
对易关系容易证明：
\begin{equation*}
[\hat{L}_i, \hat{L}^{2}] = 0\tag{C.26}
\end{equation*}
及
\begin{equation*}
[\hat{L}_i, \hat{L}_j] = \mathrm{i}\epsilon_{ijk}\hat{L}_k\tag{C.27}
\end{equation*}
交错张量（alternating tensor）$\epsilon_{ijk}$ 定义为
\begin{equation*}
\epsilon_{ijk} = \begin{cases} 1 & \text{若 } ijk \text{ 是 123 的偶置换}\\ -1 & \text{若 } ijk \text{ 是 123 的奇置换}\\ 0 & \text{若 } i, j, k \text{ 中任何两个相等} \end{cases}\tag{C.28}
\end{equation*}
式 C.27 是
\begin{equation*}
[\hat{L}_x, \hat{L}_y] = \mathrm{i}\hat{L}_z
\end{equation*}
及其循环置换的简写。该式使用了爱因斯坦求和约定：任何重复两次的指标都默认求和。
例如 $a_ib_i$ 是 $\sum_ia_ib_i$ 的简写。交错张量在如下表达式中很有用：
\begin{equation*}
(\mathbf{a}\times\mathbf{b})_i = \epsilon_{ijk}a_jb_k.
\end{equation*}
此式中 j 与 k 重复出现、默认求和。例如
\begin{equation*}
(\mathbf{a}\times\mathbf{b})_1 = \epsilon_{123}a_2b_3+\epsilon_{132}a_3b_2 = a_2b_3-a_3b_2.
\end{equation*}

算符 $\hat{L}^{2}$ 的本征函数
\begin{equation*}
|l, m_l\rangle = Y_{lm_l}(\theta,\phi) \propto P_l^{m_l}(\cos\theta)e^{\mathrm{i}m_l\phi},\tag{C.29}
\end{equation*}
称为球谐函数，本征值 $l(l+1)$；l 是角动量量子数，$P_l^{m_l}(\cos\theta)$ 是连带
勒让德多项式（见表 C.1）。

\begin{table}
\centering
\caption{$P_{l}^{|m_l|}(\cos\theta)$。}
\begin{small}
\begin{tabular}{lcccc}
\toprule
 & $m_l=0$ & $m_l=1$ & $m_l=2$ & $m_l=3$\\
\midrule
$l=0$ & 1 & -- & -- & --\\
$l=1$ & $\cos\theta$ & $\sin\theta$ & -- & --\\
$l=2$ & $\frac{1}{2}(3\cos^{2}\theta-1)$ & $\frac{1}{3}\sin\theta\cos\theta$ & $\frac{1}{3}\sin^{2}\theta$ & --\\
$l=3$ & $\frac{1}{2}(5\cos^{3}\theta-3\cos\theta)$ & $\frac{2}{3}\sin\theta(5\cos^{2}\theta-1)$ & $\frac{1}{15}\sin^{2}\theta\cos\theta$ & $\frac{1}{15}\sin^{3}\theta$\\
\bottomrule
\end{tabular}
\end{small}
\end{table}

于是
\begin{equation*}
\hat{L}^{2}|l, m_l\rangle = l(l+1)|l, m_l\rangle.\tag{C.30}
\end{equation*}
类似地
\begin{equation*}
\hat{L}_z|l, m_l\rangle = m_l|l, m_l\rangle.\tag{C.31}
\end{equation*}
升降算符 $\hat{L}_{\pm}$ 定义为
\begin{equation*}
\hat{L}_{\pm} = \hat{L}_x \pm \mathrm{i}\hat{L}_y.\tag{C.32}
\end{equation*}
由此可证
\begin{equation*}
\hat{L}_{\pm}|l, m_l\rangle = \sqrt{l(l+1)-m_l(m_l\pm1)}\,|l, m_l\pm1\rangle.\tag{C.33}
\end{equation*}

某些态没有轨道角动量。单态（S=0）就是一个好例子，原因在于它们具有实波函数——
这可以直截了当地证明。无磁场时哈密顿量为实，故若 $\psi$ 是哈密顿量 $\hat{\mathcal{H}}$
能量为 E 的本征函数，则 $\psi^{*}$ 亦然（$\hat{\mathcal{H}}\psi=E\psi$ 蕴含
$\hat{\mathcal{H}}\psi^{*}=E\psi^{*}$）。既然该态按定义是单态，$\psi=\psi^{*}$，
$\psi$ 为实。$\hat{\mathbf{L}}$ 各分量的算符都含 i，故若 $\psi$ 为实，期望值
$\langle\psi|\hat{L}_{\alpha}|\psi\rangle$（$\alpha=x,y,z$）都是纯虚数；但期望值
是可测量、必须是纯实数——因此全为零：该态没有轨道角动量。

\section{氢原子}

本书考虑的大多数原子不是氢，但薛定谔方程只能对氢原子精确求解。本节把氢原子的
一些结果列表备查。

球对称势的薛定谔方程可以用分离变量求解，本征函数 $\Psi$ 可写为径向与角向部分
之积：
\begin{equation*}
\Psi = R(r)\Theta(\theta)\Phi(\phi)\tag{C.34}
\end{equation*}
$\Phi(\phi)$ 的方程是
\begin{equation*}
\frac{\mathrm{d}^{2}\Phi}{\mathrm{d}\phi^{2}} + m_l^{2}\Phi = 0\tag{C.35}
\end{equation*}
容易解得
\begin{equation*}
\Phi = A\,\mathrm{e}^{\mathrm{i}m_l\phi}.\tag{C.36}
\end{equation*}
波函数径向部分的解可以求出，形式为
\begin{equation*}
R_{nl}(r) = -2\left(\frac{Z}{na_0}\right)^{\frac{3}{2}}\left[\frac{(n-l-1)!}{n[(n+l)!]^{3}}\right]^{\frac{1}{2}}\rho^{l}\,\mathrm{e}^{-\rho/2}L_{n+l}^{2l+1}(\rho)\tag{C.37}
\end{equation*}
n 是主量子数，
\begin{equation*}
\rho = \frac{2Zr}{na_0},\tag{C.38}
\end{equation*}
连带拉盖尔多项式为
\begin{equation*}
L_{\alpha}^{\beta}(x) = \frac{\mathrm{d}^{\beta}}{\mathrm{d}x^{\beta}}\left(\mathrm{e}^{x}\frac{\mathrm{d}^{\alpha}}{\mathrm{d}x^{\alpha}}[\mathrm{e}^{-x}x^{\alpha}]\right) = \sum_{j=0}^{\alpha-\beta}\frac{\alpha!^{2}(-1)^{j+\beta}}{(j+\beta)!j!(\alpha-\beta-j)!}x^{j}.\tag{C.39}
\end{equation*}
原子态的量子数范围是 $n\geq1$、$0\leq l\leq n-1$、$-l\leq m_l\leq l$。表 C.2 列出
n=1、2、3 的一些径向波函数。氢原子的若干径向波函数 $R(r)$ 绘于图 C.2，同时画出
径向概率密度 $r^{2}R(r)^{2}$。

\begin{table}
\centering
\caption{径向波函数。}
\begin{small}
\begin{tabular}{cccl}
\toprule
$n$ & $l$ & $L_{n+l}^{2l+1}(x)$ & $R_{nl}(\rho)$\\
\midrule
1 & 0 & $L_1^1(x)=-1$ & $2(Z/a_0)^{3/2}e^{-\rho/2}$\\
2 & 0 & $L_2^1(x)=2x-4$ & $(1/2\sqrt{2})(Z/a_0)^{3/2}(2-\rho)e^{-\rho/2}$\\
2 & 1 & $L_3^3(x)=-6$ & $(1/2\sqrt{6})(Z/a_0)^{3/2}\rho e^{-\rho/2}$\\
3 & 0 & $L_3^1(x)=-3x^{2}+18x-18$ & $(1/9\sqrt{3})(Z/a_0)^{3/2}(6-6\rho+\rho^{2})e^{-\rho/2}$\\
3 & 1 & $L_4^3(x)=24x-96$ & $(1/9\sqrt{6})(Z/a_0)^{3/2}(4-\rho)\rho e^{-\rho/2}$\\
3 & 2 & $L_5^5(x)=-120$ & $(1/9\sqrt{30})(Z/a_0)^{3/2}\rho^{2}e^{-\rho/2}$\\
\bottomrule
\end{tabular}
\end{small}
\end{table}

\begin{figure}
\centering
\includegraphics[width=0.24\textwidth]{figC_02.pdf}
\caption{n=1、2、3、4 的径向波函数 $R(r)$（左）与相应概率密度 $r^{2}R(r)^{2}$
（右）。}
\end{figure}

\section{g 因子}

磁场 B 中电子的能量为
\begin{equation*}
E = g\mu_{\mathrm{B}}m_sB\tag{C.40}
\end{equation*}
g 称为 g 因子。能级因此分裂 $g\mu_{\mathrm{B}}B$。狄拉克电子理论（超出本书范围）
的自然结论是 g 恰好等于 2。实际上 g 因子并不完全等于 2，而取值
\begin{equation*}
g = 2\left(1+\frac{\alpha}{2\pi}+\cdots\right) = 2.0023\cdots\tag{C.41}
\end{equation*}
$\alpha$ 是精细结构常数——无量纲量：
\begin{equation*}
\alpha = \frac{e^{2}}{4\pi\epsilon_0\hbar c} = \frac{1}{137.04}.\tag{C.42}
\end{equation*}
这一由量子电动力学（QED）得到的理论值与实验符合到惊人的精度。与 g=2 的偏差可以
这样解释：电磁相互作用源于虚光子的作用。电子实际上可以发射一个虚光子、随后再
吸收它。若在虚光子存续期间测量电子的磁矩，实际测到的是电子--光子对的磁矩——
包含一个额外的轨道分量。也可能在多于一个虚光子被产生的时段测量系统，不过这类
过程随虚光子数目增多而概率变小。平均而言，系统的磁矩比最初预期的略高——这解释了
$g-2$ 的非零值。该表达式是 $\alpha$ 的幂级数，逐项反映越来越曲折的虚光子产生与
吸收过程的贡献。

\section{d 轨道}

d 波函数（$l=2$）的角向部分——以凝聚态物理最常考虑的形式——可以由形如
$Y_{2m}e^{\mathrm{i}m_l\phi}$ 的函数的线性组合构造，其中 $Y_{2m}$ 是表 C.1 中
l=2 行所列的函数：
\begin{align*}
Y_{xy} &= \frac{Y_{22}-Y_{2-2}}{\mathrm{i}\sqrt{2}}\tag{C.43}\\
Y_{x^{2}-y^{2}} &= \frac{Y_{22}+Y_{2-2}}{\sqrt{2}}\tag{C.44}\\
Y_{yz} &= \frac{-Y_{21}-Y_{2-1}}{\mathrm{i}\sqrt{2}}\tag{C.45}\\
Y_{zx} &= \frac{-Y_{21}+Y_{2-1}}{\sqrt{2}}\tag{C.46}\\
Y_{z^{2}} &= Y_{20}\tag{C.47}
\end{align*}
给出如下结果（含归一化前置因子）：
\begin{align*}
Y_{xy} &= \sqrt{\frac{15}{16\pi}}\sin^{2}\theta\sin2\phi = \sqrt{\frac{15}{4\pi}}\frac{xy}{r^{2}}\tag{C.48}\\
Y_{x^{2}-y^{2}} &= \sqrt{\frac{15}{16\pi}}\sin^{2}\theta\cos2\phi = \sqrt{\frac{15}{16\pi}}\frac{x^{2}-y^{2}}{r^{2}}\tag{C.49}\\
Y_{yz} &= \sqrt{\frac{15}{4\pi}}\sin\theta\cos\theta\sin\phi = \sqrt{\frac{15}{4\pi}}\frac{yz}{r^{2}}\tag{C.50}\\
Y_{zx} &= \sqrt{\frac{15}{4\pi}}\sin\theta\cos\theta\cos\phi = \sqrt{\frac{15}{4\pi}}\frac{zx}{r^{2}}\tag{C.51}\\
Y_{z^{2}} &= \sqrt{\frac{15}{16\pi}}(3\cos^{2}\theta-1) = \sqrt{\frac{15}{16\pi}}\frac{2z^{2}-x^{2}-y^{2}}{r^{2}}.\tag{C.52}
\end{align*}

\begin{figure}
\centering
\includegraphics[width=0.35\textwidth]{figC_03.pdf}
\caption{内禀自旋--轨道相互作用。在核的参考系中电子绕核运动；但在与电子共动的
惯性系中则是核在绕行。}
\end{figure}

这些轨道连同 s、p 轨道示于图 3.1。

\section{自旋--轨道相互作用}

原子中的自旋--轨道相互作用来源如下。考虑绕原子运动的电子：图 C.3 上图描绘了核
静止参考系中的情形；下图中原子被画在与电子共动的惯性系里——核看起来在绕电子
运动。绕行的核构成电流，在原点处产生磁场
\begin{equation*}
\mathbf{B} = \frac{\boldsymbol{\varepsilon}\times\mathbf{v}}{c^{2}},\tag{C.53}
\end{equation*}
其中
\begin{equation*}
\boldsymbol{\varepsilon} = -\nabla V(r) = -\frac{\mathbf{r}}{r}\frac{\mathrm{d}V(r)}{\mathrm{d}r}\tag{C.54}
\end{equation*}
是核在电子处产生的电场，$V(r)$ 是相应的势能。式 C.53 来自狭义相对论中电磁场的
变换。这一磁场与电子自旋作用，在哈密顿量中给出
\begin{align*}
H_{\mathrm{so}} &= -\frac{1}{2}\mathbf{m}\cdot\mathbf{B}\tag{C.55}\\
&= \frac{e\hbar^{2}}{2m_{\mathrm{e}}c^{2}r}\frac{\mathrm{d}V(r)}{\mathrm{d}r}\,\mathbf{S}\cdot\mathbf{L}\tag{C.56}
\end{align*}
轨道角动量为 $\hbar\mathbf{L}=m_{\mathrm{e}}\mathbf{r}\times\mathbf{v}$，磁矩
$\mathbf{m}=(ge\hbar/2m)\mathbf{S}$；式 C.55 中的因子 $\frac{1}{2}$ 是相对论的
托马斯因子。用狄拉克方程可以极其优雅地得到这一结果——相对论修正自动包含在内。
\fn{2}{对应电子轨道上各个瞬时静止系，存在无穷多组洛伦兹变换。速度方向改变时洛伦兹变换不对易——我们的推导忽略了这一点。净洛伦兹变换包含一个转动，这一托马斯进动恰好给出二分之一的因子。}\bio{L.~H. Thomas\\（1903--1992）}这一效应称为内禀自旋--轨道相互作用，是原子中电子波函数的自旋部分与轨道部分之间的相互作用。对类氢原子中的库仑场
\begin{equation*}
\frac{1}{r}\frac{\mathrm{d}V(r)}{\mathrm{d}r} = \frac{Ze}{4\pi\epsilon_0 r^{3}}\tag{C.57}
\end{equation*}
对量子数为 l 与 n 的电子态有
\begin{equation*}
\langle r^{-3}\rangle = \frac{Z^{3}}{a_0^{3}n^{3}l(l+\frac{1}{2})(l+1)}\tag{C.58}
\end{equation*}
故自旋--轨道劈裂为
\begin{equation*}
\frac{Z^{4}e^{2}\hbar^{2}\langle\mathbf{S}\cdot\mathbf{L}\rangle}{4\pi\epsilon_0 a_0^{3}n^{3}l(l+\frac{1}{2})(l+1)}.\tag{C.59}
\end{equation*}

\section{Landé g 因子}

Landé g 因子 $g_J$ 由
\begin{equation*}
g_J = \frac{3}{2} + \frac{S(S+1)-L(L+1)}{2J(J+1)}\tag{C.60}
\end{equation*}
给出。式 C.60 中 $g_J$ 的表达式可以从磁矩算符的表达式出发得到：
\begin{equation*}
\hat{\boldsymbol{\mu}} = \mu_{\mathrm{B}}(g_L\hat{\mathbf{L}} + g_S\hat{\mathbf{S}}).\tag{C.61}
\end{equation*}
其中 $g_L=1$ 与 $g_S=2$ 分别是轨道角动量与自旋角动量的 g 因子。在许多原子中 S
与 L 可能不是好量子数，但 J 是。因此平行于 J 的磁矩分量是守恒量，垂直于 J 的
分量则不是。于是写
\begin{equation*}
\hat{\boldsymbol{\mu}} = g_J\mu_{\mathrm{B}}\hat{\mathbf{J}},\tag{C.62}
\end{equation*}
$g_J$ 是待定常数。式 C.62 意味着 $g_J$ 是 $\hat{\mathbf{L}}+2\hat{\mathbf{S}}$
在 $\hat{\mathbf{J}}$ 上的投影。把式 C.61 两边乘 $\hat{\mathbf{J}}$：
\begin{equation*}
\hat{\boldsymbol{\mu}}\cdot\mathbf{J} = \mu_{\mathrm{B}}(g_L\hat{\mathbf{L}}\cdot\hat{\mathbf{J}} + g_S\hat{\mathbf{S}}\cdot\hat{\mathbf{J}}).\tag{C.63}
\end{equation*}
把式 C.62 两边乘 $\hat{\mathbf{J}}$ 得\fn{3}{在此及后续步骤中，我们利用 $S^{2}$ 的本征值为 $S(S+1)$、$L^{2}$ 的本征值为 $L(L+1)$、$J^{2}$ 的本征值为 $J(J+1)$ 这一事实。}
\begin{equation*}
\hat{\boldsymbol{\mu}}\cdot\mathbf{J} = g_J\mu_{\mathrm{B}}\hat{\mathbf{J}}^{2} = g_J\mu_{\mathrm{B}}J(J+1).\tag{C.64}
\end{equation*}
又 $\mathbf{L}^{2}=(\mathbf{J}-\mathbf{S})^{2}=\mathbf{J}^{2}+\mathbf{S}^{2}-2\mathbf{S}\cdot\mathbf{J}$，故
\begin{equation*}
\mathbf{S}\cdot\mathbf{J} = \frac{1}{2}(\mathbf{J}^{2}-\mathbf{L}^{2}+\mathbf{S}^{2})\tag{C.65}
\end{equation*}
且 $\mathbf{S}^{2}=(\mathbf{J}-\mathbf{L})^{2}=\mathbf{J}^{2}+\mathbf{L}^{2}-2\mathbf{L}\cdot\mathbf{J}$，故
\begin{equation*}
\mathbf{L}\cdot\mathbf{J} = \frac{1}{2}(\mathbf{J}^{2}+\mathbf{L}^{2}-\mathbf{S}^{2}).\tag{C.66}
\end{equation*}
联立式 C.63 与 C.64，并代入式 C.65、C.66 的结果：
\begin{equation*}
\begin{aligned}
g_J &= g_L\left(\frac{J(J+1)+L(L+1)-S(S+1)}{2J(J+1)}\right)\\
&\quad + g_S\left(\frac{J(J+1)-L(L+1)+S(S+1)}{2J(J+1)}\right)
\end{aligned}\tag{C.67}
\end{equation*}
$g_L=1$、$g_S=2$ 时即退化为式 C.60。

\section{微扰论}

考虑哈密顿量 $\hat{\mathcal{H}}_0$，其本征函数 $|\phi_i\rangle$ 与本征值 $E_i$
已知：
\begin{equation*}
\hat{\mathcal{H}}_0|\phi_i\rangle = E_i|\phi_i\rangle.\tag{C.68}
\end{equation*}
本征函数 $|\phi_i\rangle$ 全部正交：
\begin{equation*}
\langle\phi_i|\phi_j\rangle = \int\mathrm{d}\tau\,\phi_i^{*}\phi_j = \delta_{ij}\tag{C.69}
\end{equation*}
$\mathrm{d}\tau$ 是体积元。现设在 $\hat{\mathcal{H}}_0$ 上加微扰 $\hat{V}$，新
哈密顿量为 $\hat{\mathcal{H}}$：
\begin{equation*}
\hat{\mathcal{H}} = \hat{\mathcal{H}}_0 + \hat{V}.\tag{C.70}
\end{equation*}
加微扰前系统处于能量 $E_k$ 的态 $|\phi_k\rangle$。系统的新本征函数将为 $|\psi\rangle$、
本征值为 E：
\begin{equation*}
\hat{\mathcal{H}}|\psi\rangle = E|\psi\rangle,\tag{C.71}
\end{equation*}
但 $|\psi\rangle$ 与 E 目前都未知。$|\psi\rangle$ 可以用旧本征函数展开：
\begin{equation*}
|\psi\rangle = \sum_ja_j|\phi_j\rangle.\tag{C.72}
\end{equation*}
代入式 C.71，左乘 $\langle\phi_i|$ 并在适当体积上积分：
\begin{equation*}
\sum_ja_j\int\mathrm{d}\tau\,\phi_i^{*}(\hat{\mathcal{H}}_0+\hat{V})\phi_j = E\sum_ja_j\int\mathrm{d}\tau\,\phi_i^{*}\phi_j,\tag{C.73}
\end{equation*}
用原本征函数的正交性（式 C.69）得
\begin{equation*}
\sum_jV_{ij}a_j = (E-E_i)a_i,\tag{C.74}
\end{equation*}
称为久期方程（secular equation），其中
\begin{equation*}
V_{ij} = \langle\phi_i|\hat{V}|\phi_j\rangle = \int\mathrm{d}\tau\,\phi_i^{*}\hat{V}\phi_j\tag{C.75}
\end{equation*}
是微扰的矩阵元。从现在起我们假定非简并情形，即假定没有态是简并的。

若微扰很小，$|\psi\rangle$ 将非常接近初态 $|\phi_k\rangle$：$a_k\approx1$，且
$i\neq k$ 时 $|a_i|\ll1$。既然波函数只有微小变化，可以设想新能量为
\begin{equation*}
E \approx E_k,\tag{C.76}
\end{equation*}
即完全不变。利用 $a_k\approx1$ 与 $i\neq k$ 时 $a_i\ll1$ 可以改进这一零阶近似。
把它们代入久期方程可得能量与 $a_i$ 的近似表达式。对 i=k，久期方程中的求和由
j=k 的项主导：
\begin{equation*}
E \approx E_k + V_{kk},\tag{C.77}
\end{equation*}
态能量的移动简单地就是 $V_{kk}$。这一包含“一阶微扰论”能量修正的结果在本书
正文中用得非常多。久期方程也可对 $i\neq k$ 求值：
\begin{equation*}
a_i = \frac{V_{ik}}{E_k-E_i}\tag{C.78}
\end{equation*}
（其中 E 已换成 $E_k$）。把式 C.78 代回式 C.74 得下一级能量近似：
\begin{equation*}
E \approx E_k + V_{kk} + \sum_{i\neq k}\frac{|V_{ik}|^{2}}{E_k-E_i}\tag{C.79}
\end{equation*}
这就是“二阶微扰论”结果。这一过程可以继续下去，给出 E 的幂级数表达式，级数的
逐项包含微扰矩阵元的更高次幂。

\section*{延伸阅读}

\begin{itemize}[itemsep=0.2em]
\item A.~I.~Rae,《Introduction to quantum mechanics》第 3 版, IOP Publishing
1992。
\item J.~J.~Sakurai,《Modern quantum mechanics》第 2 版, Addison-Wesley 1994。
\item P.~A.~M.~Dirac,《The principles of quantum mechanics》第 4 版, OUP 1958。
\item B.~H.~Bransden 与 C.~J.~Joachain,《Physics of atoms and molecules》,
Longman 1983。
\item P.~W.~Atkins,《Molecular quantum mechanics》, OUP 1983。
\item G.~K.~Woodgate,《Elementary atomic structure》, OUP 1980。
\end{itemize}
